\documentclass[10pt,a4paper]{amsart}
\usepackage{amsmath,amssymb,amsthm,graphicx}
\usepackage{verbatim}
\newenvironment{reference}{\comment}{\endcomment}
%\newenvironment{reference}{}{}
\newenvironment{slogan}{\comment}{\endcomment}
\newenvironment{history}{\comment}{\endcomment}

% For commutative diagrams we use Xy-pic
\usepackage[all]{xy}

% We use 2cell for 2-commutative diagrams.
\xyoption{2cell}
% Theorem environments.
%
\theoremstyle{plain}
\newtheorem{theorem}[subsection]{Theorem}
\newtheorem{proposition}[subsection]{Proposition}
\newtheorem{lemma}[subsection]{Lemma}

\theoremstyle{definition}
\newtheorem{definition}[subsection]{Definition}
\newtheorem{example}[subsection]{Example}
\newtheorem{exercise}[subsection]{Exercise}
\newtheorem{situation}[subsection]{Situation}

\theoremstyle{remark}
\newtheorem{remark}[subsection]{Remark}
\newtheorem{remarks}[subsection]{Remarks}

\numberwithin{equation}{subsection}

% Macros
%
\def\lim{\mathop{\mathrm{lim}}\nolimits}
\def\colim{\mathop{\mathrm{colim}}\nolimits}
\def\Spec{\mathop{\mathrm{Spec}}}
\def\Hom{\mathop{\mathrm{Hom}}\nolimits}
\def\Ext{\mathop{\mathrm{Ext}}\nolimits}
\def\SheafHom{\mathop{\mathcal{H}\!\mathit{om}}\nolimits}
\def\SheafExt{\mathop{\mathcal{E}\!\mathit{xt}}\nolimits}
\def\Sch{\mathit{Sch}}
\def\Mor{\mathop{\mathrm{Mor}}\nolimits}
\def\Ob{\mathop{\mathrm{Ob}}\nolimits}
\def\Sh{\mathop{\mathit{Sh}}\nolimits}
\def\NL{\mathop{N\!L}\nolimits}
\def\CH{\mathop{\mathrm{CH}}\nolimits}
\def\proetale{{pro\text{-}\acute{e}tale}}
\def\etale{{\acute{e}tale}}
\def\QCoh{\mathit{QCoh}}
\def\Ker{\mathop{\mathrm{Ker}}}
\def\Im{\mathop{\mathrm{Im}}}
\def\Coker{\mathop{\mathrm{Coker}}}
\def\Coim{\mathop{\mathrm{Coim}}}

% Boxtimes
%
\DeclareMathSymbol{\boxtimes}{\mathbin}{AMSa}{"02}

%
% Macros for moduli stacks/spaces
%
\def\QCohstack{\mathcal{QC}\!\mathit{oh}}
\def\Cohstack{\mathcal{C}\!\mathit{oh}}
\def\Spacesstack{\mathcal{S}\!\mathit{paces}}
\def\Quotfunctor{\mathrm{Quot}}
\def\Hilbfunctor{\mathrm{Hilb}}
\def\Curvesstack{\mathcal{C}\!\mathit{urves}}
\def\Polarizedstack{\mathcal{P}\!\mathit{olarized}}
\def\Complexesstack{\mathcal{C}\!\mathit{omplexes}}
% \Pic is the operator that assigns to X its picard group, usage \Pic(X)
% \Picardstack_{X/B} denotes the Picard stack of X over B
% \Picardfunctor_{X/B} denotes the Picard functor of X over B
\def\Pic{\mathop{\mathrm{Pic}}\nolimits}
\def\Picardstack{\mathcal{P}\!\mathit{ic}}
\def\Picardfunctor{\mathrm{Pic}}
\def\Deformationcategory{\mathcal{D}\!\mathit{ef}}
\usepackage{cleveref}
\crefname{originaldefinition}{original definition}{original definitions}
\Crefname{originaldefinition}{Original definition}{Original definitions}
\makeatletter
\newlabel{groupoids-definition-invariant-open}{{03BC}{source}{}{sourcepointer-1}{}}
\newlabel{groupoids-definition-invariant-open@cref}{{[originaldefinition][0][]03BC}{[1][0][]source}{}{}{}}
\makeatother
\crefname{originaldefinition}{original definition}{original definitions}
\Crefname{originaldefinition}{Original definition}{Original definitions}
\makeatletter
\newlabel{groupoids-definition-restrict-groupoid}{{02VC}{source}{}{sourcepointer-2}{}}
\newlabel{groupoids-definition-restrict-groupoid@cref}{{[originaldefinition][0][]02VC}{[1][0][]source}{}{}{}}
\makeatother
\crefname{originallemma}{original lemma}{original lemmas}
\Crefname{originallemma}{Original lemma}{Original lemmas}
\makeatletter
\newlabel{more-morphisms-lemma-separated-locally-quasi-finite-morphisms-fppf-descend}{{02W8}{source}{}{sourcepointer-3}{}}
\newlabel{more-morphisms-lemma-separated-locally-quasi-finite-morphisms-fppf-descend@cref}{{[originallemma][0][]02W8}{[1][0][]source}{}{}{}}
\makeatother
\crefname{originallemma}{original lemma}{original lemmas}
\Crefname{originallemma}{Original lemma}{Original lemmas}
\makeatletter
\newlabel{morphisms-lemma-immersion-locally-quasi-finite}{{01TN}{source}{}{sourcepointer-4}{}}
\newlabel{morphisms-lemma-immersion-locally-quasi-finite@cref}{{[originallemma][0][]01TN}{[1][0][]source}{}{}{}}
\makeatother
\crefname{originallemma}{original lemma}{original lemmas}
\Crefname{originallemma}{Original lemma}{Original lemmas}
\makeatletter
\newlabel{schemes-lemma-immersions-monomorphisms}{{01L7}{source}{}{sourcepointer-5}{}}
\newlabel{schemes-lemma-immersions-monomorphisms@cref}{{[originallemma][0][]01L7}{[1][0][]source}{}{}{}}
\makeatother
\crefname{originaldefinition}{original definition}{original definitions}
\Crefname{originaldefinition}{Original definition}{Original definitions}
\makeatletter
\newlabel{spaces-definition-immersion}{{02YU}{source}{}{sourcepointer-6}{}}
\newlabel{spaces-definition-immersion@cref}{{[originaldefinition][0][]02YU}{[1][0][]source}{}{}{}}
\makeatother
\crefname{originallemma}{original lemma}{original lemmas}
\Crefname{originallemma}{Original lemma}{Original lemmas}
\makeatletter
\newlabel{spaces-lemma-morphism-sheaves-with-P-effective-descent-etale}{{03I2}{source}{}{sourcepointer-7}{}}
\newlabel{spaces-lemma-morphism-sheaves-with-P-effective-descent-etale@cref}{{[originallemma][0][]03I2}{[1][0][]source}{}{}{}}
\makeatother
\crefname{originalsection}{original section}{original sections}
\Crefname{originalsection}{Original section}{Original sections}
\makeatletter
\newlabel{spaces-section-lists}{{02WE}{source}{}{sourcepointer-8}{}}
\newlabel{spaces-section-lists@cref}{{[originalsection][0][]02WE}{[1][0][]source}{}{}{}}
\makeatother
\crefname{originaltheorem}{original theorem}{original theorems}
\Crefname{originaltheorem}{Original theorem}{Original theorems}
\makeatletter
\newlabel{spaces-theorem-presentation}{{02WW}{source}{}{sourcepointer-9}{}}
\newlabel{spaces-theorem-presentation@cref}{{[originaltheorem][0][]02WW}{[1][0][]source}{}{}{}}
\makeatother
\title{Locally closed subspaces of an algebraic space presented by an etale equivalence relation}
\author{The Stacks Project Authors}
\date{stacks-project:spaces-properties.tex:local-snapshot-20261012; GNU FDL 1.2 or later}
\begin{document}
\maketitle
\noindent\textit{Editorial note.} All delivered mathematical text is the author's own, copied from the pinned Stacks Project chapter file spaces-properties.tex; no statement or proof step was rewritten, reordered or removed.\par\medskip
\noindent\textit{Editorial note.} The counted proof is the author's complete argument: the construction in one direction through base change along the presentation with the verification of invariance, the inverse construction by restricting the equivalence relation and applying the presentation theorem, and the closing descent of immersions along the fppf cover.\par\medskip
\noindent\textit{Editorial note.} The label lemma-subspaces-presentation names the counted statement, and the proof environment immediately following it in the source is delivered verbatim as its proof.\par\medskip
\noindent\textit{Editorial note.} The counted statement is carried by the author's own lemma environment in the chapter source.\par\medskip
\noindent\textit{Editorial note.} Pre-open gate: python operations/agent-authoring/tools/precard\_gate.py --source spaces-properties.tex --label lemma-subspaces-presentation --start 1631 --end 1681 --local-id 28044 returned RESULT=OK (L28044 free, occupancy FREE, public interval clear). Candidate taken from sources/stacks-spaces-properties/candidates-halfA.json.\par\medskip
\noindent\textit{Editorial note.} The statement and proof are copied from the pinned chapter file; the cited inputs are kept as external pointers with their canonical Stacks tags.\par\medskip
\noindent\textit{Allowed standard prerequisites (editorial).}\begin{itemize}
\item Presentation of an algebraic space as the quotient of an etale equivalence relation. Spaces, the theorem presenting an algebraic space as the quotient of an etale equivalence relation on a scheme. Supplies the identification X = U/R and the resulting description of subsheaves via restriction.
\item Restriction of an equivalence relation to an invariant subscheme. Groupoids, the definition of an invariant open and of the restriction of a groupoid to an invariant subscheme. Supplies the equivalence relation R' on Z' and hence the algebraic space Z = Z'/R'.
\item Descent of immersions for the fppf topology. More on Morphisms, the lemma that separated and locally quasi-finite morphisms descend along fppf coverings, together with the facts that immersions are monomorphisms, separated and locally quasi-finite. Supplies the descent of the property immersion from the cover U to the space X.
\end{itemize}\medskip
\ifdefined\hypertarget\hypertarget{sourcepointer-1}{}\fi
\noindent\textit{Source reference.} Original definition 03BC: sources/stacks-chow/tags.txt, source lines 4343--4343. Prerequisite pointer invoked at proof line(s) 1651; it names a standard fact of a sibling chapter that the counted argument uses, and every step of the counted proof built on it is retained verbatim in the delivered body.\par\medskip
\ifdefined\hypertarget\hypertarget{sourcepointer-2}{}\fi
\noindent\textit{Source reference.} Original definition 02VC: sources/stacks-chow/tags.txt, source lines 3763--3763. Prerequisite pointer invoked at proof line(s) 1656; it names a standard fact of a sibling chapter that the counted argument uses, and every step of the counted proof built on it is retained verbatim in the delivered body.\par\medskip
\ifdefined\hypertarget\hypertarget{sourcepointer-3}{}\fi
\noindent\textit{Source reference.} Original lemma 02W8: sources/stacks-chow/tags.txt, source lines 3795--3795. Prerequisite pointer invoked at proof line(s) 1669; it names a standard fact of a sibling chapter that the counted argument uses, and every step of the counted proof built on it is retained verbatim in the delivered body.\par\medskip
\ifdefined\hypertarget\hypertarget{sourcepointer-4}{}\fi
\noindent\textit{Source reference.} Original lemma 01TN: sources/stacks-chow/tags.txt, source lines 2380--2380. Prerequisite pointer invoked at proof line(s) 1667; it names a standard fact of a sibling chapter that the counted argument uses, and every step of the counted proof built on it is retained verbatim in the delivered body.\par\medskip
\ifdefined\hypertarget\hypertarget{sourcepointer-5}{}\fi
\noindent\textit{Source reference.} Original lemma 01L7: sources/stacks-chow/tags.txt, source lines 2076--2076. Prerequisite pointer invoked at proof line(s) 1665; it names a standard fact of a sibling chapter that the counted argument uses, and every step of the counted proof built on it is retained verbatim in the delivered body.\par\medskip
\ifdefined\hypertarget\hypertarget{sourcepointer-6}{}\fi
\noindent\textit{Source reference.} Original definition 02YU: sources/stacks-chow/tags.txt, source lines 3889--3889. Prerequisite pointer invoked at proof line(s) 1677; it names a standard fact of a sibling chapter that the counted argument uses, and every step of the counted proof built on it is retained verbatim in the delivered body.\par\medskip
\ifdefined\hypertarget\hypertarget{sourcepointer-7}{}\fi
\noindent\textit{Source reference.} Original lemma 03I2: sources/stacks-chow/tags.txt, source lines 4588--4588. Prerequisite pointer invoked at proof line(s) 1672; it names a standard fact of a sibling chapter that the counted argument uses, and every step of the counted proof built on it is retained verbatim in the delivered body.\par\medskip
\ifdefined\hypertarget\hypertarget{sourcepointer-8}{}\fi
\noindent\textit{Source reference.} Original section 02WE: sources/stacks-chow/tags.txt, source lines 3801--3801. Prerequisite pointer invoked at proof line(s) 1663; it names a standard fact of a sibling chapter that the counted argument uses, and every step of the counted proof built on it is retained verbatim in the delivered body.\par\medskip
\ifdefined\hypertarget\hypertarget{sourcepointer-9}{}\fi
\noindent\textit{Source reference.} Original theorem 02WW: sources/stacks-chow/tags.txt, source lines 3819--3819. Prerequisite pointer invoked at proof line(s) 1635; it names a standard fact of a sibling chapter that the counted argument uses, and every step of the counted proof built on it is retained verbatim in the delivered body.\par\medskip
\begin{lemma}
\label{lemma-subspaces-presentation}
Let $S$ be a scheme. Let $j : R \to U \times_S U$ be an \'etale equivalence
relation. Let $X = U/R$ be the associated algebraic space
(Spaces, Theorem \ref{spaces-theorem-presentation}). There is a
canonical bijection
$$
R\text{-invariant locally closed subschemes }Z'\text{ of }U
\leftrightarrow
\text{locally closed subspaces }Z\text{ of }X
$$
Moreover, if $Z \to X$ is closed (resp.\ open) if and only if
$Z' \to U$ is closed (resp.\ open).
\end{lemma}

\begin{proof}
Denote $\varphi : U \to X$ the canonical map. The bijection sends
$Z \to X$ to $Z' = Z \times_X U \to U$. It is immediate from the definition
that $Z' \to U$ is an immersion, resp.\ closed immersion, resp.\ open
immersion if $Z \to X$ is so. It is also clear that $Z'$ is $R$-invariant
(see Groupoids, Definition \ref{groupoids-definition-invariant-open}).

\medskip\noindent
Conversely, assume that $Z' \to U$ is an immersion which is $R$-invariant.
Let $R'$ be the restriction of $R$ to $Z'$, see
Groupoids, Definition \ref{groupoids-definition-restrict-groupoid}.
Since $R' = R \times_{s, U} Z' = Z' \times_{U, t} R$ in this case
we see that $R'$ is an \'etale equivalence relation on $Z'$. By
Spaces, Theorem \ref{spaces-theorem-presentation} we see
$Z = Z'/R'$ is an algebraic space. By construction we have
$U \times_X Z = Z'$, so $U \times_X Z \to Z$ is an immersion.
Note that the property ``immersion'' is preserved under base change
and fppf local on the base (see Spaces, Section \ref{spaces-section-lists}).
Moreover, immersions are separated and locally quasi-finite (see
Schemes, Lemma \ref{schemes-lemma-immersions-monomorphisms}
and
Morphisms, Lemma \ref{morphisms-lemma-immersion-locally-quasi-finite}).
Hence by More on Morphisms, Lemma
\ref{more-morphisms-lemma-separated-locally-quasi-finite-morphisms-fppf-descend}
immersions satisfy descent for fppf covering. This means all the hypotheses of
Spaces,
Lemma \ref{spaces-lemma-morphism-sheaves-with-P-effective-descent-etale}
are satisfied for $Z \to X$, $\mathcal{P}=$``immersion'',
and the \'etale surjective morphism $U \to X$. We conclude that $Z \to X$
is representable and an immersion, which is the
definition of a subspace (see
Spaces, Definition \ref{spaces-definition-immersion}).

\medskip\noindent
It is clear that these constructions are inverse to each other and we win.
\end{proof}

\end{document}
